\section{Conclusion}
\label{sec:conclusion}

This study asked whether a day-ahead electricity-price forecasting advantage survives explicit information constraints and translates into economic value when forecasting, uncertainty and trading rules are frozen before final evaluation.

Using public ENTSO-E data for the Germany--Luxembourg market, models were developed on 2019--2025 data and January--July 2026 was reserved for a first-look holdout. The final design distinguished a Full information set from a Tier-1 robustness set, used chronological model selection, constructed uncertainty from delivery-day-safe out-of-sample residuals, and mapped forecasts into a pre-specified single-cycle storage decision rule.

The primary result survives the holdout. Full XGBoost attains an MAE of €17.51/MWh compared with €29.33/MWh for lag-24. Under the primary economic assumptions, the Full point strategy earns €20,002.02, or €1,437.60 more than the lag-24 strategy. The forecast-value contrast remains positive in all nine pre-specified efficiency--cost cells, satisfying the mechanical holdout confirmation criterion.

The secondary results qualify that conclusion in useful ways. Tier-1 XGBoost remains more accurate and economically valuable than lag-24, but it is materially weaker than Full, showing that the unresolved renewable-vintage issue is economically consequential. The uncertainty-aware Full strategy sacrifices €872.01 of cumulative net P\&L but raises the trade hit rate to 99.4\% and reduces the worst daily outcome to -€0.83. Uncertainty therefore acts as a conservative abstention mechanism rather than a profit-maximizing layer under the tested rule.

The paper's contribution is not a new forecasting algorithm and not the observation that forecast accuracy and economic value can diverge. Instead, it demonstrates a reproducible forecast-to-decision workflow in which information timing, model selection, uncertainty construction, economic assumptions and the final evaluation are treated as one research design.

The principal unresolved issue is now clear. Before the Full model can be interpreted as a genuinely deployable 11:45 D-1 system, historical or live renewable forecast vintages must be verified at that exact decision origin. Future work should also move directly to 15-minute forecasts and decisions, represent multivariate price-path uncertainty, include more complete battery and market-access costs, and accumulate a longer post-transition holdout for more credible tail-risk inference.

The broader lesson is that useful forecast evaluation begins after the forecast is produced. A forecast is operationally valuable only if the information used was genuinely available, the decision rule was specified before outcomes were known, and the claimed economic gain survives assumptions that matter to the decision.
